
\emph{Mache quantitative Angaben und veranschauliche diese mit Diagrammen. 
\\
Geben  im Abschnitt Ergebnisse keine zusätzlichen Informationen zu deinen Methoden an (z. B. die Methode zur Auswahl der Parameter oder das Versuchsprotokoll). 
\\
Geben hier nur die reinen Fakten an; interpretiere deine Ergebnisse noch nicht. Das heißt, keine Sätze wie: "Die verbesserte Leistung zeigt, dass der Controller effektiv ist..." oder "Ein wahrscheinlicher Grund für diese Beobachtung ist ...".}



%---------------------------------------------------------
\subsection{Die schriftliche Ausarbeitung}
\label{sec:schriftlich}




\subsubsection{Umfang}
Die schriftliche Ausarbeitung sollte so kurz wie möglich und so umfangreich wie nötig sein. Bitte sprich das mit dem/der Betreuer/in ab.

Die schriftliche Abschlussarbeit soll eine geordnete Darstellung der im Verlauf der Arbeit angestellten Überlegungen, durchgeführten Ableitungen, entwickelten Algorithmen, Experimenten usw. im Sinne einer Diskussion bringen. Das Herausarbeiten und Bewerten der Ergebnisse ist dabei die wesentliche geistige Leistung. Eine konzentrierte Darstellung (etwa $30$ bis $50$ A4- Seiten) ist anzustreben; vermeide unbedingt weitschweifige Abhandlungen. Experimentprotokolle, Rechnerausdrucke, etc. können geordnet, eindeutig und vollständig beschriftet als Anhang beigeheftet wer-den. Berechnungen sind trotz der gebotenen Kürze so ausführlich zu bringen, sodass eine schrittweise Überprüfung möglich ist.

\subsubsection{Äußere Form}

Die Arbeit ist in Deutsch oder Englisch abzufassen und vorzugsweise mit dem Textverarbeitungssystem \LaTeX  zu schreiben. Die bereitgestellte Vorlage kann genutzt werden.
\hyperlink{https://www.overleaf.com/project/629f45ce28c12f2c468b8f78}{
Vorlage für Abschlussarbeit an der HTW}

\begin{itemize}
\item  Verwende niemals fremde Ideen/ Texte/ Bilder, ohne die Urheberrechte zu nennen und zu klären. 
\item  Beginne mit einer logischen Struktur, exakte Formulierungen mache später. Mehr dazu in Kap.~\ref{sec:struktur}
\item Lass dir deine Arbeit von jemandem überarbeiten.
\item  Die Zusammenfassung sollte besonders gut werden. Viele Leser werden nur diese lesen.
\item  Auch wenn für die endgültige Fassung kein Inhaltsverzeichnis benötigt wird, macht es Sinn, mit Hilfe des Befehls 
\texttt{$\backslash$tableofcontents} während des Schreibens eines zu erstellen. Damit behälst du deine Struktur im Auge. % Use the command \tableofcontents}for that purpose. 
\item Bemühe dich um Symmetrie in Bezug auf die Anzahl und Länge der Abschnitte pro Kapitel oder der Unterabschnitte pro Abschnitt (Einleitung/ Schluss sind Ausnahmen mit weniger Unterüberschriften). 



\item Zwischen Abschnittsüberschriften und Unterabschnittsüberschriften sollten keine "losen" Textabschnitte eingefügt werden. Ein oder zwei Sätze sind möglich, aber nur für einen Überblick über den gesamten Abschnitt.


\item Ein Abschnitt kann nicht nur einen einzigen Unterabschnitt enthalten.

\item Überspringe keine Überschriftenebenen.


\item Vermeide mehr als drei Ebenen (Abschnitt, Unterabschnitt, Unterunterabschnitt). Wenn du auf einer tieferen Ebene unterteilen musst, überarbeite deine Struktur.

\item Jeder Absatz sollte eine eindeutige Idee/ Botschaft enthalten. Außerdem sollte der Lesende das Thema des Absatzes bereits in den ersten drei Wörtern, mindestens aber im ersten Satz erkennen.

\item Achte auf eine logische Verbindung innerhalb und zwischen den Absätzen. Am einfachsten ist es, zuerst die Hauptgedanken aufzuschreiben und sie erst dann zu Absätzen auszubauen.

\item  Werde nicht kreativ und verwende verschiedene Zeilenumbrüche zwischen Abschnitten. Wenn das passiert, denke noch einmal genau über deine Dokumentstruktur nach.

\item Der oder die Lesende sollte schnell wissen, wo er oder sie welche Informationen finden kann. Deshalb ist wichtig, klare Überschriften zu verwenden. Benutze keine Abkürzungen in Überschriften.

\end{itemize}
%---------------------------------------------------------
\subsection{Kolloquium}
\label{sec:kolloquium}

Als Abschluss der Abschlussarbeit findet ein Kolloquium statt. Der Termin wird rechtzeitig bekannt gegeben. 

Beim Kolloquium sind in der Regel die Gutachter:innen, Mitarbeiter:innen, Studierende  und andere interessierte Gäste zugegen. Du referierst zunächst in einem möglichst frei zu haltenden Kurzvortrag von nicht mehr als 15~Minuten Dauer über wesentliche Ergebnisse deiner Arbeit. Es folgt eine Diskussion und Fragerunde von maximal 30- 45 Minuten. Im Anschluss an den Seminarvortrag ist nach Möglichkeit eine Vorführung der konkreten Ergebnisse erwünscht.


\subsubsection{Ablauf}

 Die Präsentation sollte für einen Beamer gestaltet und ca. 15 min lang sein.
 
Es folgt die eigentliche Verteidigung, in der die Prüfungskommision dich zu deiner Arbeit befragt. Zum Abschluss berät sich die Prüfungskommission  unter Ausschluss der Öffentlichkeit über die zu vergebende Note . Du erfährst das Ergebnis der Notenfindung unmittelbar nach der Beratung. Sofern du bestanden hast, bekommst du sofort eine Bescheinigung über den erfolgreichen Abschluss deines Studiums. 

\subsubsection{Folien}

\begin{itemize}
\item 
Für den Vortrag ist eine erneute Auswahl und Bewertung des in der schriftlichen Fassung der Abschlussarbeit niedergelegten Stoffes notwendig. Durch den Kurzvortrag soll beim Zuhörerkreis, der im allgemeinen mit der Problematik deiner Arbeit nicht vertrautes Fachpublikum ist, Verständnis und Interesse für die Arbeit wecken und erzielte Ergebnisse erläutern. Eine didaktisch klare und straffe Gliederung deiner Überlegungen ist unerlässlich. 

\item
Gebe dem Vortrag eine klare Struktur Er beginnt normalerweise mit einer kurzen Inhaltsübersicht, auf die eine knappe Einleitung folgt, in der du grundsätzlich in die Thematik deiner Abschlussarbeit einführst. Im Hauptteil der Präsentation beschreibst du die Resultate deiner Arbeit, wobei du davon ausgehen darfst, dass dein Publikum fachkundig und mit dem Thema vertraut ist. Trotzdem solltest du Begriffe, soweit diese nicht zum allgemeinen Sprachgebrauch  fachvertrauter Spezialist:innen zählen, kurz erläutern. Zum Ende des Vortrags empfiehlt sich ein Ausblick auf mögliche Verbesserungen, Erweiterungen usw.

    \item Folienvorlagen findest du in der HTW Cloud \hyperlink{
https://corporatedesign.htw-berlin.de/musterdokumente/powerpoint-vorlagen/}{\emph{Folienvorlage HTW}} 
\item Latex in ppt
\hyperlink{http://www.jonathanleroux.org/software/iguanatex/}{latexpluin}
\item Verwende einen gut lesbaren und großen Schrifttyp. 
\item Vermeide den Einsatz von mehr als zwei unterschiedlichen Fonts auf einer Folie.
\item 
Verzichte auf Effekte. 
\item Bei Bildern und auch bei Tabellen solltest du darauf achten, diese nicht zu überladen. Arbeite jeweils die Informationen heraus, die für den Vortrag relevant sind. Das geht sehr einfach mit selbst gestalteten Bildern. Bilder sollten selbstverständlich den Rand einer Folie nicht überragen. 
\item Problematisch ist die Verwendung von Screenshots. Durch die geringe Auflösung der meisten Beamer wirken solche Bilder schnell „pixelig“. 
\item Bei allen Sachverhalten, die du nicht selbst erarbeitet hast, musst du Quellen benennen. Das gilt auch für deine Folien. Für die Folien, die Quellen immer dort anzugeben, wo sie zitiert werden – z.B. in Kleinschrift direkt unter einem Bild. 
\item Zur Abschätzung der Anzahl der benötigten Folien gehe bitte davon aus, dass pro Folie etwa zwei Minnten Vortragszeit anzusetzen sind.
 \item Bevorzuge auf den Folien Stichpunkte statt ausgeschriebene Sätze. Dadurch vermeidest du, dass dein Publikum Texte liest, statt dir zuzuhören. 
% \item Handouts?
% \item Backup Folien
\end{itemize}

\subsubsection{Präsentation} 

\begin{itemize}

\item 
Stelle sicher, dass am Tage des Abschlusskolloquiums keine technischen Probleme auftreten. So sind nicht in allen Räumen der HTW Beamer mit einem HDMI Eingang vorhanden. Auch solltest du wissen, wie dein Laptop zu konfigurieren ist, um die Folien über den Beamer auszugeben und wie der Bildschirmschoner ausgeschaltet werden kann, damit er nicht mitten im Vortrag aktiviert wird. Falls du eine Vorführung planst, musst du dich  darum kümmern, dass dafür alles notwendige bereit steht und auch lauffähig ist. Sofern das Abschlusskolloquium in der HTW stattfindet, steht i.Allg. ein Whiteboard zur Verfügung, das due zur Beantwortung von Fragen nutzen kannst. Bereite dich auf die Präsentation gut vor. 

 %Eine Faustregel ist maximal vier Worte pro Zeile und sechs Zeilen pro Folien zu verwenden (Ausnahmen bestätigen die Regel). 

\item Versuche, den Vortrag frei zu halten. 
\item Mache dir Kärtchen, auf denen Stichpunkte stehen.
\item 
Rede laut und akzentuiert. 
\item Treten während deines Vortrags Fragen auf, so musst du abwägen, wie du verfahren willst. Es gibt Fragen, die sofort beantwortet werden müssen, weil sie essentiell für das Verständnis deines Vortrags sind. In den allermeisten Fällen kannst du  die Fragen aber auch an das Ende deines Vortrags verschieben.
\item Ein Probevortrag vor deinem Betreuer/ deiner Betreuerin wird dringend nahegelegt.



%\item Wissenschaftlicher Vortrag!!
\end{itemize}


%---------------------------------------------------------

\subsection{Datenablage}
\label{sec:datenablage}

Für die Archivierung der Arbeit ist eine fundierte Datenablage
%in der HTW Could 
erwünscht. Sie soll enthalten:

\begin{enumerate}
\item alle Quellen, die zur Erstellung der Berichte benutzt worden sind (Das sind also zum Beispiel: .tex-File, .bib-File, Stylefiles, alle Graphiken, bitte mit absoluten/ relativen Pfaden aufpassen, dass alles ohne Fehler kompilierbar ist...),


\item die gehaltenen Präsentationen, wieder mit allen benötigten Quellen wie Graphiken, .tex-Files, ppt-File...,
\item alle Daten, die zum Reproduzieren der Ergebnisse benötigt werden. (z.B. Simulationen, die im Bericht gezeigten Daten, Daten zu Messkurven)
enthält.
\end{enumerate}
%---------------------------------------------------------
